\section*{零基础阅读指南（先看我！）}
\addcontentsline{toc}{section}{零基础阅读指南}

\paragraph{这门课到底在学什么？}一句话：\textbf{用计算机"做实验"来解决统计问题}。
传统数理统计靠公式推导，这门课靠"让计算机模拟成千上万次随机实验，
然后数数/求平均"。整本资料的主线只有一条：

\begin{center}
\fbox{\parbox{0.92\textwidth}{\centering
第1--2章：教计算机"掷骰子"（造出服从各种分布的随机数）\\[1mm]
第3--4章：用这些"骰子"做实验，检验统计方法好不好用\\[1mm]
第5章：让实验更省钱（同样次数，误差更小）\\[1mm]
第6章：手里只有一份数据、没有总体时怎么办（Bootstrap）\\[1mm]
第7章：数据缺了一块时怎么估参数（EM 算法）}}
\end{center}

\paragraph{怎么用这份资料？}
\begin{itemize}[leftmargin=2em]
  \item 每章先读开头的"故事"，知道这章为什么存在；
  \item 灰色\textbf{"人话解读"}框：把公式翻译成大白话，看不懂公式就先看它；
  \item 蓝框=要背的公式定理，绿框=算法步骤，红框=易错点，橙框=例题；
  \item 时间紧：直接背每章末尾"考点地图"+做完最后的模拟卷。
\end{itemize}

\paragraph{常用符号速查表。}
\begin{center}
\begin{tabular}{ll}
$U\sim U(0,1)$ & $0$ 到 $1$ 之间的均匀随机数（计算机的"原料"）\\
$\E X$ & $X$ 的期望（平均值）\qquad $\Var X$：方差（波动大小）\\
$\Cov(X,Y),\ \rho$ & 协方差/相关系数（两个量一起波动的程度）\\
$F(x)$ / $p(x)$ & 分布函数（累积概率）/ 密度函数\\
$X\sim N(\mu,\sigma^2)$ & 正态分布；$\Phi$ 为标准正态分布函数\\
$\hat\theta$ & 参数 $\theta$ 的估计量（戴帽子=用数据算出来的）\\
$I(\cdot)$ & 指示函数：括号里的事发生取 $1$，否则取 $0$\\
$iid$ & 独立同分布（每个数据独立、来自同一总体）\\
$u_{q},\ t_q(n)$ & 标准正态 / $t$ 分布的 $q$ 分位数（查表得到）\\
\end{tabular}
\end{center}

\newpage
